% Styl PJATK — odtwarza wygląd z pjatk-template.typ:
% EB Garamond 10pt, justowany tekst, wytłuszczone rozdziały, numeracja 1.1,
% listingi z numerami po prawej.

% --- Główna czcionka: EB Garamond z lokalnych plików TTF -----------------
% Główna czcionka: EB Garamond (Regular / Italic / Bold / BoldItalic).
% Dodatkowe wagi (Medium, SemiBold, ExtraBold) dostępne są w contents/assets/fonts/
% i można je podłączyć przez \newfontfamily, jeśli będą potrzebne w konkretnym miejscu.
\setmainfont{EBGaramond}[
  Path            = contents/assets/fonts/,
  Extension       = .ttf,
  UprightFont     = *-Regular,
  ItalicFont      = *-Italic,
  BoldFont        = *-Bold,
  BoldItalicFont  = *-BoldItalic
]

% --- Czcionka monospace (dla listingów, \texttt, \code) -------------------
\setmonofont[Scale=MatchLowercase]{Menlo}

% --- Czcionka bezszeryfowa (opcjonalna) -----------------------------------
\setsansfont[Scale=MatchLowercase]{Helvetica}

% Rozmiar bazowy (10pt — tak jak Typst). Użyj \small/\large w razie potrzeby.
\renewcommand{\normalsize}{\fontsize{10pt}{14pt}\selectfont}
\normalsize

% --- Numeracja sekcji 1.1 i głębokość TOC --------------------------------
\setcounter{secnumdepth}{3}
\setcounter{tocdepth}{2}

% --- Stylizacja rozdziałów (titlesec) ------------------------------------
\titleformat{\chapter}[display]
  {\normalfont\Huge\bfseries}
  {\chaptertitlename\ \thechapter}
  {20pt}
  {\Huge}
\titlespacing*{\chapter}{0pt}{-30pt}{20pt}

\titleformat{\section}
  {\normalfont\Large\bfseries}{\thesection}{1em}{}
\titleformat{\subsection}
  {\normalfont\large\bfseries}{\thesubsection}{1em}{}
\titleformat{\subsubsection}
  {\normalfont\normalsize\bfseries\itshape}{\thesubsubsection}{1em}{}

% Nagłówki bez przenoszenia wyrazów (jak w Typst)
\pretolerance=2000
\tolerance=3000

% --- Nagłówki i stopki (fancyhdr) ----------------------------------------
\pagestyle{fancy}
\fancyhf{}
\fancyhead[LE]{\small\nouppercase{\leftmark}}
\fancyhead[RO]{\small\nouppercase{\rightmark}}
\fancyfoot[LE,RO]{\thepage}
\renewcommand{\headrulewidth}{0.4pt}
\renewcommand{\footrulewidth}{0pt}

\fancypagestyle{plain}{%
  \fancyhf{}
  \fancyfoot[LE,RO]{\thepage}
  \renewcommand{\headrulewidth}{0pt}
}

% --- Inline code: \code{...} z jasnoniebieskim tłem ----------------------
\newcommand{\code}[1]{%
  \colorbox{codebg}{\strut\kern0.1em\texttt{\small #1}\kern0.1em}%
}

% --- Style listingów (numeracja po prawej, brak ramki, łamanie linii) ----
\lstdefinestyle{pjatk}{%
  basicstyle        = \ttfamily\footnotesize,
  keywordstyle      = \color{codeKeyword}\bfseries,
  commentstyle      = \color{codeComment}\itshape,
  stringstyle       = \color{codeString},
  numbers           = right,
  numberstyle       = \tiny\color{codeNumGray},
  numbersep         = 8pt,
  breaklines        = true,
  breakatwhitespace = true,
  frame             = none,
  captionpos        = t,
  showstringspaces  = false,
  tabsize           = 2,
  columns           = fullflexible,
  keepspaces        = true,
  xleftmargin       = 0pt,
  xrightmargin      = 0pt,
  aboveskip         = 0.8\baselineskip,
  belowskip         = 0.8\baselineskip
}
\lstset{style=pjatk}

% Polskie nazwy list dla pakietu listings
\renewcommand{\lstlistlistingname}{Spis listingów}
\renewcommand{\lstlistingname}{Listing}

% --- Drobiazgi typograficzne ---------------------------------------------
\frenchspacing
\clubpenalty=10000
\widowpenalty=10000
\displaywidowpenalty=10000
